% SPDX-License-Identifier: Apache-2.0
\section{A Field the Netlist Names Differently}
\label{sec:x-rename}

A unit's field is a signal of its netlist, so its name has to be a
name the netlist can carry. VHDL reserves \code{next} and
\code{buffer}, Verilog reserves \code{signed} and \code{buf}, and a
field named for one of those writes a netlist that will not analyse,
so the derive refuses such a name where the field is declared. That
refusal is right and it is also a cost: the name that reads best in
Rust is sometimes exactly the one a target reserves.

\code{\#[rename("...")]} on a field says what the netlist should call
it. Rust keeps the name the code is written with; the netlist takes
the other, everywhere, in its declarations, its drives and the trace
the run records. The queue here calls its two registers \code{next}
and \code{buffer}, which is what they are, and the netlist calls them
\code{nxt} and \code{held}.

The lowering reads the \code{impl} and never the struct, so it writes
every reference under the name Rust knows and the rename is applied
to the lowered unit afterwards, from a table the derive writes beside
the field names. The check for a reserved word then runs on the
netlist's name rather than on Rust's, and two fields that would take
one name in the netlist are refused, since the netlist would declare
it twice.

\begin{figure}[htbp]
\centering
\input{rename_timing.tex}
\caption{The queue: \code{nxt} is the word about to leave and
\code{held} the one behind it, under the names the netlist gives
them.}
\label{fig:x-rename}
\end{figure}

\lstinputlisting[caption={\code{ex\_rename.rs}. Run with \code{bazel run //lib/examples:ex\_rename}.},label={lst:x-rename}]{lib/examples/ex_rename.rs}

\lstinputlisting[style=ascii,caption={What \code{ex\_rename} printed when the build ran it: the queue cycle by cycle, then the Verilog, with \code{nxt} and \code{held} in it.},label={lst:x-rename-out}]{out_rename.txt}
